\documentclass[a5paper,twoside]{article}
\usepackage[utf8]{inputenc}
\usepackage{lmodern}
\usepackage[T1]{fontenc}
\usepackage[hungarian]{babel}

\usepackage{graphicx}
\graphicspath{ {images/} }

\usepackage[bookmarks]{hyperref}

\usepackage[chorded]{songs}
% \usepackage[lyric]{songs}

\usepackage{draftwatermark}
\SetWatermarkText{3. Minta}
\SetWatermarkHorCenter{7cm}
\SetWatermarkScale{0.8}
\SetWatermarkLightness{0.9}

\title{Budaörsi Antiochia dalgyűjtemény akkordokkal}
\author{Balics Eszter - Donkó István}
% \date{2017. február 5.}

% \title{Budaörsi Antiochia daloskönyv}
% \date{2017. február 5.}
% \author{Készítette Donkó István - Benkő Matyi munkájából kiindulva}

\setlength{\topmargin}{-1.5cm}

\setlength{\headheight}{0cm}
\setlength{\headsep}{0cm}

\setlength{\oddsidemargin}{-1.3cm}
\setlength{\evensidemargin}{-2cm}

\setlength{\textwidth}{13cm}
\setlength{\textheight}{18cm}

\setlength{\songnumwidth}{1cm}
\setlength{\versenumwidth}{0.5cm}
% \setlength{\cbarwidth}{0.1cm}

\catcode`_=12 % use underscore as regular character
\renewcommand{\_}[1]{\underline{#1}} % shorter command for adding melismas
\songcolumns{1}
\notenames{A}{H}{C}{D}{E}{F}{G}

% \renewcommand{\lyricfont}[1]{\bf#1}
% \renewcommand{\lyricfont}[1]{\rmfamily\bf#1}

% \usepackage[outline]{contour}
% \contourlength{0.1pt}
% \contournumber{10}
% \renewcommand{\lyricfont}[1][asdf]{\contour{black}{#1}}

\newindex{titleidx}{antiochia_title_index}

\begin{document}

  \begin{titlepage}
    \setlength{\oddsidemargin}{-1.3cm}

    \vspace*{1.3cm}
    \includegraphics[width=11cm]{Logo_songbook}

    \thispagestyle{empty}
  \end{titlepage}

  \renewcommand{\idxtitlefont}{\sffamily}
  \renewcommand{\idxlyricfont}{\rmfamily\slshape}
  \showindex{Tartalomjegyzék}{titleidx}

  \versesep=12pt plus 3pt minus 5pt

  % \songsection{Antiochiás énekek}
  \begin{songs}{titleidx}
    \beginsong{Add tovább}
[
  index={Nagy tűz, hatalmas láng}
]

  \ifchorded
    \beginverse*
      {\nolyrics Előjáték: \[D] \[F#m] \[G] \[A] \[A7]}
    \endverse
  \fi

  \beginverse\memorize
    Nagy \[D]tűz, hatalmas \[F#m]láng, csak egy \[G]szikra az, amitől \[A]kigyullad. \echo{amitől \[A7]kigyullad}
    \[D]Melegét szórja \[F#m]rád, csak a \[G]kezed kell, hogy oda\[A]ny{\_{ú}}jtsad. \echo{hogy oda\[A7]ny{\_{ú}}jtsad}
  \endverse

  \beginchorus
    \[G]Épp így az Isten \[D]szeretetét, \[Em] ha egyszer el\[A]fogad\[D]tad,
    Már \[G]nem kí\[D]vánsz \[G]semmi \[D]mást, csak \[G]azt, hogy \[A]add to\[D]vább.
  \endchorus

  \beginverse
    Az ^első napsu^gár, mi a ^dermedt tavaszi ^földre száll, \echo{tavaszi ^földre száll}
    ^Melegét szórja ^rád, és a ^fák rügyeiket ^kibontják. \echo{^kibontják}
  \endverse

  \beginverse*
    Refrén
  \endverse

  \beginverse
    ^Szívemben boldog^ság, mert az ^Úr, az Isten az ^én Atyám. \echo{igen, az ^én Atyám}
    ^Lelkemben ég a ^vágy, h\_{o}gy ^megtegyem az Ő ^akaratát. \echo{az Ő ^akaratát}
  \endverse

  \beginverse*
    /: A \[G]hegy csúcsáról \[D]kiáltok: Uram! \[Em] Tárd fel a \[A]vilá\[D]gom!
    \[G]Hogy to\[D]vább \[G]adjam \[D]én a \[G]szeretet \[A]Iste\[D]nét. :/ \rep{2}

    % /: A \[G]hegy csúcsáról \[D]kiáltok: Uram! \[Em]Tárd fel a \[A]vilá\[D]gom!
    % Hogy \[G]tovább \[Em]adjam \[A]én a \[G]szeretet \[A]Iste\[D]nét. :/ \rep{2}
  \endverse

\endsong

\beginsong{Adjunk hálát}[]

  \versesep=12pt plus 3pt minus 8pt

  \ifchorded
    \beginverse*
      {\nolyrics Előjáték: \[G] \[C] \[G]}
    \endverse
  \fi

  \beginverse\memorize
    Adjunk \[C]h\_{á}\[G]lát az \[C]\_{Ú}r\[G]nak,
    Aki \[Em]meghalt \[Am]és fel\[D]támadt,
    Ő a \[C]mennybe \[G]ment és a \[C]trónon \[G]ül,
    Onnan \[Em]várjuk, \[Am]míg el\[D]jő.
  \endverse

  \beginchorus
    Allelu\[Hm]ja, az Úr az \[Em]Isten!
    Allelu\[Am]ja, a Megvál\[D]tó!
    Allelu\[Hm]ja, az Úr az \[Em]Isten!
    Zengj alle\[Am]luját! \[D7]{\_{Á}}\[G]men! \[C] \[G]
  \endchorus

  \beginverse
    Az Ő ^t\_{e}s^te és ^v\_{é}^re
    Táplál ^minket ^itt a ^Földön,
    Gyere ^és ö^rülj, velünk ^ünne^pelj,
    Tiszta ^szívvel ^áldjuk ^Őt!
  \endverse

  \beginverse*
    Refrén
  \endverse

  \beginverse
    Aki ^benne ^hisz, sose ^szomja^zik,
    Ő az ^élő ^víz fo^rrása.
    Szent^l\_{é}lek, ^jöjj és ^szállj le ^ránk,
    Tölts be ^minket, ^Szeretet! ^
  \endverse

  \beginverse*
    Refrén
  \endverse

\endsong

\input{songs/aldasoddal-megyunk.tex}
\beginsong{Áldjad, én lelkem, az Urat}[]

  \ifchorded
    \beginverse*
      {\nolyrics Előjáték: \[Am] \[Em] \[Am] \[Em] \[Am]}
    \endverse
  \fi

  \beginverse\memorize
    \[Am]Áldjad, én \[Em]lelkem, az U\[Am]rat! \[Em]
    \[Am]Áld\[Em]jad, én \[C]lelkem, \[D7]az U\[E]rat! \echo{Alle\[E7]luj{\_{a}}}
    \[Am]Áldjad, én \[Em]{l\_{e}}lkem,
    \[F]Áldjad, én \[G6]{l\_{e}}l\[C]kem,
    \[Am]Áldjad, én \[Em]lelkem, az U\[Am]rat! \[Em] \[Am] \[Em]
  \endverse

  \beginverse
    ^Zengj alle^luját az Úr^nak! ^
    ^Zengj ^alle^luját ^az Úr^nak! \echo{Alle^luj{\_{a}}}
    ^Zengj alle^{l\_{u}}ját,
    ^Zengj alle^{l\_{u}}^ját,
    ^Zengj alle^luját az Úr^nak! ^ ^ ^
  \endverse

  \beginverse
    ^Dicsérjük ^együtt az U^rat! ^
    ^Di^csérjük ^együtt ^az U^rat! \echo{Alle^luj{\_{a}}}
    ^Dicsérjük ^{\_{e}}gyütt,
    ^Dicsérjük ^{\_{e}}gy^ütt,
    ^Dicsérjük ^együtt az U^rat! ^ ^ ^
  \endverse

  \beginverse
    ^Áldjad, én ^bensőm, az U^rat! ^
    ^Áld^jad, én ^bensőm, ^az U^rat! \echo{Alle^luj{\_{a}}}
    ^Áldjad, én ^{b\_{e}}nsőm,
    ^Áldjad, én ^{b\_{e}}n^sőm,
    ^Áldjad, én ^bensőm, az U^rat! ^ ^ ^
  \endverse

  \ifchorded
    \beginverse*
      {\nolyrics Utójáték: \[Am] \[Em] \[Am] \[Em] \[Am]}
    \endverse
  \fi

\endsong

\beginsong{Áldjátok az Urat}[]

  \beginverse*
    /: \[Em]Áldjátok az Ur\[Am]at, áldjátok szent ne\[Em]vét,
    Kik \[Em]házában álltok \[D]napról n\_{a}pra, \[Hm]dicsérjétek szüntelen \[Em]Őt! :/ \rep{2}
  \endverse

  \beginverse*
    /: \[Em]Tárjátok kezei\[Am]tek az él\_{ő} Isten fe\[Em]lé!
    \[Am]Áldjon meg téged \[Em]Sionból az Úr, \[D]ragyogtassa \[Hm]arcát re\[Em]ád! :/ \rep{2}
  \endverse

\endsong

\beginsong{Áldott légy, Uram}
[
  by={Taizé / Jacques Berthier}
]

  \beginverse*
    \[Dm]Áldott \[G]légy, U\[Dm]ram, szent \[A#]neved \[C]áldja \[F]lelkem! \[A]
    \[Dm]Áldott \[G]légy, U\[Dm]ram, mert \[A#]megvál\[C]tottál \[Dm]már.
  \endverse

\endsong

\beginsong{Alleluja}
[
  index={Krisztus az Úr}
]

  \versesep=12pt plus 3pt minus 8pt

  \ifchorded
    \beginverse*
      {\nolyrics Előjáték: \[D] \[F#m] \[G] \[D] \[G] \[D] \[A7] \[D]}
    \endverse
  \fi

  \beginchorus
    \[D]Allelu\[F#m]ja, alle\[G]lu – allelu\[D]ja,
    \[G]Allelu\[D]ja, alle\[A]lu\[A7]ja.
    \[D]Allelu\[F#m]ja, alle\[G]lu – allelu\[D]ja,
    \[G]Allelu\[D]ja, alle\[A]l\[A7]u\[D]ja.
  \endchorus

  \beginverse\memorize
    \[D]Krisztus az \[F#m]Úr, erő s \[G]hatalom Ö\[D]vé,
    \[G]Bátran bízz \[D]benne, hisz’ \[A]mond\[A7]ta:
    ``Ki \[D]újra ko\[F#m]pog, annak \[G]ajtót nyitnak \[D]majd''.
    \[G]Allelu\[D]ja, alle\[A]l\[A7]u\[D]ja.
  \endverse

  \beginverse*
    Refrén
  \endverse

  \beginverse
    ^Föltámadt ^Krisztus, le^győzte a ha^lált,
    ^Mindenre ^új hajnal ^vir^rad.
    ^Indulj az ^úton és ^soha meg ne ^állj!
    ^Allelu^ja, alle^l^u^ja!
  \endverse

  \beginverse*
    Refrén
  \endverse

  \beginverse
    ^Szívünkbe ^add, Uram, ^csönded, örö^möd,
    ^Víg kedvünk ^hadd legyen ^jó^ra!
    ^Bízva ki ^él, annak ^ünnep ada^tik.
    ^Allelu^ja, alle^l^u^ja!
  \endverse

  \beginverse*
    Refrén
  \endverse

  \beginverse
    ^Küzdeni ^kell, napról ^napra elkí^sér
    ^Szívünkben ^száz fényes ^ó^ra.
    ^Bízva ki ^él, annak ^öröm ada^tik.
    ^Allelu^ja, alle^l^u^ja!
  \endverse

  \beginverse*
    Refrén
  \endverse

\endsong

\beginsong{Álmaimból, lelkemből}
[
  li={Kék könyv / 12.B}
]

  \ifchorded
    \beginverse*
       {\nolyrics Előjáték:
       \[Hm] \[F#m] \[G] \[F#m] \[G] \[G] \[A] \[A]
       \[Hm] \[F#m] \[G] \[F#m] \[G] \[A] \[A] \[G] \[G] \[D] \[A]}
    \endverse
  \fi

  \beginverse\memorize
    \[D] Szürke a \[G]föld és szürke az \[D]ég, kopár az \[G]út és csillag sem \[D]ég,
    Elcsüggedt \[Em]szívű, magányos \[A]vándor, hova \[D]m\[A]\_{é}sz?
    \[D] Ne menj to\[G]vább, nem futhatsz \[D]el, csak rajtad \[G]áll, csak tenned \[D]kell,
    Elcsüggedt \[Em]szívű, magányos \[A]vándor, van re\[D]m\[A]\_{é}ny!
  \endverse

  \beginchorus
    /: \[A] Álmaim\[Hm]ból, lelkem\[F#m]ből, két kezem\[G]ből, szívem\[F#m]ből,
    Szól ez a \[G]dal, tied a \[A]dal,
    Zengjen a \[Hm]dal szíved\[F#m]ben, gyulladjon \[G]fény a lelked\[F#m]ben,
    Vezesse \[G]lépteidet az \[A]Isten, vezesse \[G]életed! \[G] \[D] :/ \rep{2}
  \endchorus

  \beginverse
    ^ Valami^kor így éltem ^én, azt hittem ^már nincs is re^mény
    Elcsüggedt ^szívű, magányos ^vándor voltam {^^\_{é}}n!
    ^ Valaki ^jött, rám neve^tett: Ne menj to^vább, nyújtsd a ke^zed,
    Elcsüggedt ^szívű, magányos ^vándor, sose ^f^\_{é}lj!
  \endverse

  \beginverse*
    Refrén
  \endverse

\endsong

\input{songs/antioch-dal.tex}
\beginsong{Apokrif ima - !!}[]

  \ifchorded
    \beginverse*
      {\nolyrics Előjáték: \[D] \[D#] \[Em] \[A6] \[A7]}
    \endverse
  \fi

  \beginverse\memorize
    \[D]Oldani Uram, s oldódni \[A]v\_{á}gyom,
    \[A7]Üdvözíteni, és üdvözülni \[D]v\_{á}gyom,
    \[G] Alkotni \[D]Uram, \[A7] s alakulni \[D]vágyom,
    \[Em] Alkotni \[A6]Uram, \[A7] s alakulni \[D]vágyom.
  \endverse

  \beginverse
    ^Ékesíteni, és ékesedni ^v\_{á}gyom,
    ^Dalolni Uram, és dallá lenni ^v\_{á}gyom,
    ^ Követeddé ^lenni, ^ Érted tenni ^vágyom,
    ^ Követeddé ^lenni, ^ Érted tenni ^vágyom.
  \endverse

  \beginverse
    ^Lámpád vagyok, ha l\_{á}tsz \_{e}n^gem,
    ^Ajtód vagyok, ha zörgetsz ^b\_{e}nnem,
    ^ Ki látod ^mit teszek, ^ hallgass el ^engem,
    ^ Ki látod ^mit teszek, ^ hallgass el ^engem.
  \endverse

\endsong

\beginsong{Ároni áldás}[]

  \ifchorded
    \beginverse*
      {\nolyrics Előjáték: \[C] \[Em] \[Dm] \[G] \[C] \[G] \[C]}
    \endverse
  \fi

  \beginverse*\memorize
    \[C]Áldjon m\_{e}g \[Em]téged az {\[Dm G]{\_{Ú}}}r!
    \[C]Áldjon m\_{e}g \[Em]téged az {\[Dm G]{\_{Ú}}}r!
    És \[C]őrizzen \[G]meg tége\[C]det! \[C7]
  \endverse

  \beginverse*
    /: \[Dm]Világo\[Em]sítsa meg az \[C]{\_{Ú}}r
    Az \[Dm]Ő or\[Em]cáját te\[C]rajtad,
    És \[F]könyörüljön teraj\[G]tad! \[G7]
    \[C]Fordítsa az \[Em]Úr az Ő \[Dm]orcáját \[G]rád!
    \[C]Fordítsa az \[Em]Úr az Ő \[Dm]orcáját \[G]rád!
    És \[C]adjon bé\[G]két tené\[C]ked! :/ \rep{2} \echo{És \[C]adjon bé\[G]két tené\[C]ked!}
  \endverse

\endsong

\input{songs/atya-teremto.tex}
\beginsong{Atyám, két kezedben}[]

  \beginverse\memorize
    \[Dm]Atyám, két ke\[C]zedben, \[A7]csak ott lakha\[Dm]tom,
    \[F]Biztonságot \[C]csak Tőled ka\[Dm]pok. \[G]
    \[Dm]Újjá így te\[C]remtesz, \[A7]sebem ápo\[Dm]lod,
    \[F]Boldogság, hogy \[C]itt van ottho\[Dm]nom. \[G]
  \endverse

  \beginchorus
    /: \[F]Tarts meg két ke\[C]zedben, \[A7]őrizz meg, U\[Dm]ram,
    \[F]Oltalmadban \[C]rejtsd el sorso\[Dm]mat! \[G] :/ \rep{2}
  \endchorus

  \beginverse
    ^Atyám, két ke^zedben ^teljes az ö^röm,
    ^Ajándékod^ban gyönyörkö^döm! ^
    ^Tékozlóként ^éltem, ^tárva most ka^rod,
    ^Hűtlenségem ^nem hánytorga^tod! ^
  \endverse

  \beginverse*
    Refrén
  \endverse

  \beginverse
    ^Atyám, két ke^zedben ^bátran sírha^tok,
    ^Fájdalmaim ^hordozod, ^tudom. ^
    ^Ott fenn a ke^reszten ^áldó két ke^zed
    ^Bűneimmel ^én szegeztem ^fel! ^
  \endverse

  \beginverse*
    Refrén \ifchorded{\nolyrics (A legvégén \[G] nélkül.)}\fi
  \endverse

\endsong

\beginsong{Atyám, Tied vagyok}[]

  \beginverse\memorize
    \[C]Atyám, Tied \[F]vagyok, és \[G]szüntelen i\[C]mádlak!
    \[C]Dicsőség \[F]szent neved\[G4]nek! \[G7]
  \endverse

  \beginchorus
    \[C]Dicsőség ne\[F]vednek, \[E]dicsőség ne\[Am]vednek,
    \[F]Dicsős\_{é}g \[G]szent neved\[F]nek! \[C]
  \endchorus

  \beginverse
    ^Jézus, Tied ^vagyok, és ^szüntelen i^mádlak!
    ^Dicsőség ^szent neved^nek! ^
  \endverse

  \beginverse*
    Refrén
  \endverse

  \beginverse
    ^Lélek, Tied ^vagyok, és ^szüntelen i^mádlak!
    ^Dicsőség ^szent neved^nek! ^
  \endverse

  \beginverse*
    Refrén \rep{2}
  \endverse

\endsong

\beginsong{Ave Maria}[]

  \beginchorus
    \[G]A\[D]ve Ma\[C]ri\[D]a, \[G]gra\[Am]tia \[C]ple\[D]na,
    \[G]D{\_{o}}mi\[Am]nus \[D]t{\_{e}}\[G]cum \[C]b{\_{e}}ne\[D]d{\_{i}}cta \[G]tu.
  \endchorus

  \beginverse\memorize
    \[G]Jézus, \[D]nékünk \[C]adtál a\[D]nyát,
    Ő \[G]kér és \[Hm]áld, ha \[C]szívünk \[D]fáj.
    \[G]Hála Jé\[Am]zus, \[D]hála Né\[G]ked,
    Mert \[C]van egy \[D]drága i\[G]mánk!
  \endverse

  \beginverse*
    Refrén
  \endverse

  \beginverse
    ^Térden ^állva ^hívunk Té^ged,
    ^Vidd el ^álmunk ^Krisztus e^lé!
    ^Jézus any^ja, ^hű Mári^a,
    Ó ^add, hogy ^lássuk Fi^ad!
  \endverse

  \beginverse*
    Refrén
  \endverse

\endsong

\beginsong{A béke napja közel}[]
[
  by={Borka Zsolt}
]

  \beginverse*
    A \[F]béke \[C]napja köz{\[Dm A#]e}l, \[F]nézd, már \[A#]fénylik a \[C]jel! \[C7]
    \[F]Jézus \[A]szíve a \[A#]földön, az \[F]é\[C]gen \[F]mindent \[C]egybeöl\[F]el.
  \endverse

\endsong

\beginsong{Bízom Benned}[]

  \ifchorded
    \beginverse*
      {\nolyrics Előjáték: \[G] \[C9] \[Em7] \[D]}
    \endverse
  \fi

  \beginchorus
    \[G]Bízom \[C9]Benned, \[Em7]igazságos \[D]Úr vagy,
    \[G]Bízom \[C9]Benned, \[Em7]irgalmas, \[D]j\_{ó}!
    \[G]Bízom \[C9]Benned \[Em7]mind\_{e}n\[D]nap, örö\[G]kké! \[C9] \[Em7] \[D]
  \endchorus

  \beginverse\memorize
    \[G] Olyan messziről jöttem \[Am7]én,
    És Te \[F]megvár\[C]tál, sokat \[G] vártál \[D]r\_{á}m,
    \[G] Úgy hívlak a magány éjje\[Am7]lén,
    Mikor \[F]elvész a \[C]szó, mikor csak \[G]átölelni \[D]jó!
    \[Am]Tudom, hogy \[Em]távol voltam, \[G]fájt Ne\[D]ked.
    \[Am]Tudom, hogy \[Em]nélküled nem \[C]élhe\[D]t\_{e}k.
  \endverse

  \beginverse*
    Refrén
  \endverse

  \beginverse
    ^ Nézd, a szívem már Ti^éd!
    Olyan ^boldog^ság, mikor ^ Téged ^l\_{á}t!
    ^ Hallod hangom lágy ne^szét,
    Kérlek, ^szólj hoz^zám, mikor csak ^Téged hív a ^szám!
    ^Életem ^nem féltem, hisz ^jó Ve^led.
    ^Tudom, hogy ^nélküled nem ^élhe^t\_{e}k!
  \endverse

  \beginverse*
    Refrén \rep{2}
  \endverse

\endsong

\input{songs/boldogsag-orom.tex}
\beginsong{Csak Jézusnak szolgálok}[]

  \beginchorus
    Csak \[G]Jézus\[D]nak \[Em]szolg\_{á}\[C]lok, az \[G]Ő or\[C]szágát \[G]v\_{á}\[D]rom,
    Mert \[G]eljön \[D]Ő, az \[C]Üdvözí\[D#]tő, ó \[G]Jézus, \[D]feltámadt Ki\[G]rály! \[C] \[G]
  \endchorus

  \beginverse\memorize
    \[Em]Áldjad, lelkem \[Hm]Őt, \[Em]dicsérd Terem\[Hm]tőd, \[G]most és \[C]mindörökké, \[G]\_{á}\[D]men!
    \[Em]Áldjad, minden \[Hm]nép, \[C]dicsérd szent Ne\[G]vét, \[G]most és \[C]mindörökké, \[D]\_{á}\[G]men!
  \endverse

  \beginverse
    ^Hálát adja^tok, ^áldást mondja^tok, ^most és ^mindörökké, ^\_{á}^men!
    ^Egyházad di^csér ^nagy jóságo^dért, ^most és ^mindörökké, ^\_{á}^men!
  \endverse

  \beginverse
    ^Erős vár az ^Úr, ^menedéket ^nyújt, ^most és ^mindörökké, ^\_{á}^men!
    ^Áldott irga^lom, ^véred olta^lom, ^most és ^mindörökké, ^\_{á}^men!
  \endverse

  \beginverse
    ^Hatalmas az ^Úr, ^jóságos az ^Úr, ^most és ^mindörökké, ^\_{á}^men!
    ^Szívem áldja ^Őt, ^ujjongjon a ^Föld, ^most és ^mindörökké, ^\_{á}^men!
  \endverse

  \beginverse
    ^Szentek, angya^lok, ^szeráfkóru^sok, ^most és ^mindörökké, ^\_{á}^men!
    ^Néked zenge^nek, ^így ünnepel^nek, ^most és ^mindörökké, ^\_{á}^men!
  \endverse

  \beginverse
    ^Áldott légy ^Atya, ^s egyszülött ^Fia, ^most és ^mindörökké, ^\_{á}^men!
    ^Éltető Lé^lek, ^Téged dicsér^lek, ^most és ^mindörökké, ^\_{á}^men!
  \endverse

\endsong

\beginsong{Ébredj ember}
[
  li={https://youtu.be/-DGzHCfmv5k alapján.}
]

  \versesep=12pt plus 3pt minus 8pt

  \ifchorded
    \beginverse*
      {\nolyrics Előjáték: \[C] \[F] \[C] \[C] \[F] \[G] \[Am] \[F] \[C] \[Em] \[Am] \[Dm] \[G] \[C]}
    \endverse
  \fi

  \beginverse
    Ébredj, \[C]ember, \[G] mély álmod\[Am]ból, \[F] Jézus \[C]megment \[Dm] rabságod\[G]ból!
    Hogyha \[Am]hívod, \[Em] eljön \[F]Ő, egész \[C]bensőm, \[Dm] imádjad \[G]Őt!
  \endverse

  \beginchorus
    Alle\[C]luja, alle\[F]l\_{u}\[C]ja!
    Alle\[C]luja, alle\[F]l\_{u}\[G]ja!
    Alle\[Am]luja, alle\[F]luja, al\[C]le\[Em]l\_{u}\[Am]ja!
    Alle\[Dm]luja, alle\[G]l\_{u}\[C]ja!
  \endchorus

  \beginverse*
    {\nolyrics \echo{Két félhang emelkedés.} \[A7]}
  \endverse

%   \textnote{két félhang emelkedés}

  \beginchorus
%     {\nolyrics (két félhang emelkedés) \[A7]}
    Alle\[D]luja, alle\[G]l\_{u}\[D]ja!
    Alle\[F#m]luja, alle\[G]l\_{u}\[A]ja!
    Alle\[Hm]luja, alle\[G]luja, al\[D]le\[F#]l\_{u}\[Hm]ja!
    Alle\[Em]luja, alle\[A]l\_{u}\[D]ja!
  \endchorus

  \beginverse\memorize
    Áldjad, \[D]lelkem, \[A] az Úr ne\[Hm]vét, \[G] mert meg\[D]tartja \[Em] ígére\[A]tét!
    Mennyi \[Hm]jót tett \[F#m] az Úr ve\[G]lem, Reá \[D]bízom \[Em] az éle\[A]tem.
  \endverse

%   \beginverse
%     Föltámadt Jézus, győztes király, elhozta nékünk az országát.
%     Jöjj Szentlélek, szállj közénk, hozz a mennyből szép, tiszta fényt!
%   \endverse

  \beginverse*
    Refrén
  \endverse

  \beginverse
    És az ^élet, ^íme, megjele^nt, ^ köztünk ^él Ő ^ Szentlelké^ben.
    Hív Ő ^téged, ^ hogy élj ve^le, add át ^néki ^ a szíve^det!
  \endverse

  \beginverse*
    Refrén
  \endverse

\endsong

\beginsong{Együtt fog ujjongni}[]

  \beginverse*
    \[H7]Együtt fog ujjongni \[Em]fiatal és öreg is
    És \[Am]örömében táncot lejt a \[H7]l\_{á}ny.
  \endverse

  \beginverse*
    % {\nolyrics  \[H7] \[Em] \[Am] \[H7]}

    \[H7]La, la la la, la la \[Em]la, la-la-la la la la
    La \[Am]la la la la, la la la la, \[H7]la la la la
  \endverse

  \beginverse*
    \[H7]Gyászukat ünnepre fordítom, s nem búsulnak ezután!
    Vi\[Am]dámmá te\[D]szem a szívü\[G]ket,
    Hogy \[Em]örvendezze\[Am]nek a \[H7]bánatuk he\[Em]ly\_{e}tt!
    Vi\[Am]dámmá te\[D]szem a szívü\[G]ket,
    Hogy \[Em]örvendezze\[Am]nek a \[H7]bánatuk he\[Em]lyett!
  \endverse

\endsong

\beginsong{Én Uram, én Istenem}[]

  \ifchorded
    \beginverse*
      {\nolyrics Előjáték: \[Em] \[H7] \[G] \[D] \[Em] \[H7] \[G] \[D] \[Em]}
    \endverse
  \fi

  \beginverse\memorize
    \[Em]Én Uram, én \[H7]Istenem,
    \[G]Vedd el tőlem \[D]mindenem,
    Ami \[Am]gátol \[Am7] Fe\[Em]léd!
  \endverse

  \beginverse
    ^Én Uram, én ^Istenem,
    ^Add meg nekem ^mindenem,
    Ami ^segít ^ Fe^léd!
  \endverse

  \beginverse
    \[Em]Én Uram, én \[H7]Istenem,
    \[G]Fogadd el az \[D]életem,
    \[C]Hadd legyen e\[D]gészen a Ti{\[G H7]{\_{é}}}d!
  \endverse

  \beginverse
    \[Em]Én Uram, én \[H7]Istenem,
    \[G]Fogadd el az \[D]életem,
    \[C]Hadd legyen e\[D]gészen a Ti\[Em]éd! \[H7]
  \endverse

  \ifchorded
    \beginverse*
      {\nolyrics \[G] \[D] \[Am] \[Am7] \[Em]}
    \endverse
  \fi

\endsong


\input{songs/feny-ami-bennem-eg.tex}
\beginsong{Föltámadt! Alleluja!}
[
  by={Sillye Jenő},
  index={Álltok szótlanul}
]

  % \ifchorded
  %   \beginverse*
  %     {\nolyrics Előjáték: \[E] \[A] \[A7]}
  %   \endverse
  % \fi

  \beginverse\memorize
    \[D]Álltok \[A7]sz{\_{ó}}tla\[D]nul, a \[G]dárda \[A]földre \[D]hull, \[Hm]
    Itt \[G]más e\[A]rő az {\[D]\[Hm]{\_{Ú}}}r, itt \[Em]más e\[A]r{\_{ő}} az \[D]Úr!
  \endverse

  \beginverse
    Az ^éj az ^h{\_{o}}sszú ^volt, de ^eltűnt ^már a ^Hold, ^
    A ^hajnal {^felsik}{^^{\_{o}}}lt, a ^hajnal ^f{\_{e}}lsi^kolt: \[D7]
  \endverse

  \beginchorus
    \[G]Föltá\[A]madt! Allel{\[D]\[Hm]{\_{u}}}ja! \[Em]Föltámadt! \[A7]Allel\_{u}\[D]ja! \[Hm]
    \[G]Föltá\[A]madt! Allel{\[D]\[Hm]{\_{u}}}ja! \[Em]Föltámadt! \[A7]Allel\_{u}\[D]ja!
  \endchorus

  \beginverse
    ^Elég volt ^m{\_{á}}r, e^lég, ^fénnyel ^hinti az ^ég ^
    Az ^új nap {^reggel}{^^{\_{é}}}t, az ^új nap ^r{\_{e}}gge^lét.
  \endverse

  \beginverse
    ^Virágban ^{\_{á}}ll a ^rét, a ^virágos ^úton ^át ^
    ^Valaki ^jön fel{^^{\_{é}}}d, ^valaki ^j{\_{ö}}n fe^léd. \[D7]
  \endverse

  \beginverse*
    Refrén
  \endverse

  \beginverse
    ^Virágban ^{\_{á}}ll a ^rét, a ^virágos ^úton ^át ^
    ^Jézus ^jön fel{^^{\_{é}}}d, ^Jézus ^j{\_{ö}}n fe^léd.
  \endverse

  \beginverse*
    Refrén \rep{2}
  \endverse

\endsong

\input{songs/gyertek-dicserjuk-egyutt.tex}
\input{songs/ha-jon-az-ur.tex}
\beginsong{Hálát adok az esti órán}
[
  ititle={Esti hála}
]

  % / jelek után: https://www.youtube.com/watch?v=ggrTwNKP5Ew

  \beginverse\memorize
    \[D]Hálá\[Hm]t adok az \[Em]esti \[A6]órán,
    \[D]Hálá\[D7]t adok, hogy \[G]itt az \[A7]éj,
    \[D]Hálá\[D7]t adok, hogy \[G]szívem \[Gm]mélyén
    \[D]Hála \[A7]dala \[D]kél.
  \endverse

  \beginverse
    ^Hálá^t adok a ^csillag^fényért,
    ^Hálá^t a sűrű ^éjje^lért,
    ^Hálá^t adok, hogy ^nem vagy ^messze,
    ^Most is ^küldesz ^fényt.
  \endverse

  \beginverse
    ^Hálá^t adok, hogy ^küldtél ^testvért
    ^Társkén^t, aki ma ^mellém ^állt.
    ^Hálá^t adok, hogy ^megen^gedted
    ^Azt is, ^ami ^fájt.
  \endverse

  \beginverse
    ^Hálá^t adok, hogy ^Tested ^táplált,
    % Igéd
    ^Hálá^t, hogy ma is ^szólt sza^vad.
    ^Hálá^t adok, hogy ^lelkem ^biztos
    ^Cél fe^lé ha^lad.
  \endverse

  \beginverse
    ^Hálá^t adok, hogy ^mellém ^álltál,
    ^Hálá^t, hogy szép az ^ottho^nom.
    % van jó?
    ^Hálá^t, hogy munkám ^elvé^geztem,
    ^S békén ^alha^tom.
  \endverse

  \beginverse
    ^Hálá^t adok a ^sok-sok ^jóért,
    ^Hálá^t adok, hogy ^megbo^csátsz.
    ^Hálá^t adok, hogy ^el nem ^hagytál,
    ^S újabb ^útra ^vársz.
  \endverse

  \beginverse
    ^Hálá^t adok, hogy ^szemed ^rajtam,
    ^Hálá^t adok, hogy ^reményt ^adsz.
    ^Hálá^t adok, hogy ^biztos ^eljössz,
    ^Mint a ^virra^dat.
  \endverse

\endsong

\beginsong{Hálát adok, hogy itt a reggel}
[
  ititle={Reggeli hála}
]

  \beginverse\memorize
    \[D]Hálá\[Hm]t adok, hogy \[Em]itt a \[A6]reggel,
    \[D]Hálá\[D7]t adok az \[G]új na\[A7]pon.
    \[D]Hálá\[D7]t adok, hogy \[G]minden \[Gm]percem
    \[D]Néked \[A7]adha\[D]tom.
  \endverse

  \beginverse
    ^Hálá^t nem csak a ^jó test^vérért,
    ^Hálá^t mindenki^ért a^dok.
    ^Hálá^t adok, hogy ^minden ^sértést
    ^Megbo^csátha^tok.
  \endverse

  \beginverse
    ^Hálá^t adok, hogy ^munkát ^küldesz,
    ^Hálá^t adok, hogy ^fény ra^gyog.
    ^Hálá^t adok, hogy ^lelkem ^fényes
    ^És bol^dog va^gyok.
  \endverse

  \beginverse
    ^Hálá^t adok a ^vidám ^percért,
    ^Hálá^t, ha szomo^rú le^szek.
    ^Hálá^t adok, hogy ^adsz e^lém célt
    ^És ke^zed ve^zet.
  \endverse

  \beginverse
    ^Hálá^t adok, hogy ^hallom ^hangod,
    ^Hálá^t adok a ^jó hí^rért.
    ^Hálá^t adok, hogy ^Tested ^adtad
    ^Minden ^embe^rért.
  \endverse

  \beginverse
    ^Hálá^t adok az ^üdvös^ségért,
    ^Hálá^t adok, hogy ^várni ^fogsz.
    ^Hálá^t adok, hogy ^ma reg^gel még
    ^Hálát ^adha^tok.
  \endverse

\endsong

\beginsong{Hazamegy a vándor}
[
  by={Pintér Béla}
]

  \beginverse
    \[H] Ami ezen a földön \[C#m]körbevesz, \[F#] az olyan nekem, mint egy \[H]múló álom.
    \[G#m] Nemrég kezdődött, máris \[E]vége lesz, ha a \[F#]nap felkel a láthatáron.
    \[H] Nincs itt maradandó \[C#m]városom, \[F#] a szívem azt az \[H]új világot várja.
    \[G#m] És összehajtom majd a \[E]vánkosom, ha a \[F#]hajnal rám talál.
  \endverse

  \beginchorus
    Hazamegy a \[H]vándor, erről álmodott ő \[G#m] százszor.
    Arra tette fel az \[E]életét, hogy egy \[G#m]szép napon majd \[F#sus4]célba \_{é}r. \[F#]
    Hazamegy a \[H]vándor, eltűnik, mint a \[G#m] kámfor.
    Mintha soha nem lett \[E]volna itt, egy szempillantás és meg\[H]érkez\_{i}k. \[F#]
  \endchorus

  \beginverse*
    1. versszak
  \endverse

  \beginverse*
    Refrén
  \endverse

  \beginverse*
    Mert \[E]új eget, és egy új \[G#m]f\_{ö}ldet \[F#]v\_{á}runk,
    \[E]Új folyókban úszunk majd, és \[G#m]új hegyeket \[F#]m\_{á}szunk.
    \[E] Ez a bolygó nem az \[G#m]igazi vi\[F#]l\_{á}gunk.
    Itt csak \[E]utazunk, át\[F#]utazunk.
  \endverse

  \beginverse*
    Refrén \rep{2} \[H]
  \endverse

\endsong

\beginsong{Hús-vér templom}
[
  by={Pintér Béla}
]

  \ifchorded
    \beginverse*
      {\nolyrics Előjáték: /: \[Dm] \[A#] \[C] \[F] :/ \rep{2}}
      % TODO: A# helyett Gm?
    \endverse
  \fi

  \beginchorus
    /: Nekem \[Dm]nincs más \[Gm]Rajtad \[C]kívül, \[F]Jézus,
    \[Dm]Éle\[Gm]tem for\[C]rása \[F]vagy.
    \[Dm]Kézzel, \[Gm7]lábbal, \[A]szívvel, \[Dm]szájjal
    \[A#]Dicsérlek, Ura\[A4]m. \[A] :/ \rep{2}
  \endchorus

  \beginverse
    Ez a \[Dm]hús-vér \[Gm]templom \[A]Érted \[Dm]épült,
    Neked \[A#]ég a \[Gm7]tűz bent, \[A]oltárai\[Dm]nál.
    Ez a hús-vér \[Gm]templom \[A]Téged \[Dm]dicsér,
    Szívem \[A#]minden \[Gm7]húrja \[A]Rólad muzsi\[Dm]kál.
  \endverse

  \beginverse*
    Refrén
  \endverse

\endsong

\input{songs/im-itt-jon-o.tex}
\beginsong{Ímé, az Isten sátora}[]

  \ifchorded
    \beginverse*
      {\nolyrics Előjáték: \[Em] \[D] \[Em] \[Hm] \[Em]}
    \endverse
  \fi

  \beginverse\memorize
    \[Em]Ímé, az \_{I}sten \[Hm]sátora az \[Em]emberekkel \[Hm]van,
    \[Em]Maga az \_{I}sten \[D]lesz velük, az \[G]ő Iste\[H]nük. \[H7]
  \endverse

  \beginchorus
    /: \[Em]J\_{é}zus, Te \[D]vagy a Király!
    \[Em]J\_{é}zus, Te \[D]vagy a Ki\[Em]rály! :/ \rep{2} \[Hm]
  \endchorus

  \beginverse
    ^Isten letörli a ^könnyeket az ^ő szemeik^ről.
    ^Halál nem lesz már ^t\_{ö}bbé, sem ^gyász, sem fájda^lom. ^
  \endverse

  \beginverse*
    Refrén
  \endverse

  \beginverse
    Ki ^győz, örökségül ^nyer mindent, szí^vébe írom ne^vem,
    ^Annak \_{I}ste^ne leszek, s fi^am lesz ő né^kem. ^
  \endverse

  \beginverse*
    Refrén \rep{2}
  \endverse

\endsong

\input{songs/indulj-a-vilagba.tex}
\input{songs/indulj-es-menj.tex}
\beginsong{Itt vagyok most, jó Uram}[]

  \beginverse\memorize
    \[D]Itt vagyok most, jó Uram, \[G]Hozzád száll e dal.
    \[D]Tudom, sokszor hívtalak, de ma \[G]másképp szól e hang,
    Mert ma \[Em]nem azt kérem, hogy \[A]adj még, nem kö\[D]nyörgök: figyelj r\[G]ám,
    Csak \[D]azt mon\[Hm7]dom, hogy \[Em7]szeret\[A7]lek, Aty\[D]ám! \[D7]
  \endverse

  \beginchorus
    Szeret\[Em7]lek, ó, U\[A7]ram, hálám \[D]szívemből fak\[G]ad,
    Hiszen \[Em7]mindent látsz Te ben\[A]nem, \[D]hatalmad olyan \[D7]nagy!
    Szeret\[Em7]lek, ó, U\[A]ram, hálám \[D]szívemből fak\[G]ad,
    És \[D]ének\[Hm7]lem, hogy \[Em7]áldott \[A7]légy, U\[D]ram!
  \endchorus

  \beginverse
    ^Ismét térden állok én, s ^tudom, hogy hallgatsz rám.
    Bár ^sokszor hideg szavakat szól ^kopott és halk imám,
    De most ^el kell, hogy mondjam egy^szer, amit ^oly rég szeretn^ék,
    Hogy ^szeretl^ek, a ^szívem ^csak Ti^éd! ^
  \endverse

  \beginverse*
    Refrén
  \endverse

\endsong

\beginsong{Jézus a mi oltalmunk}[]

  \versesep=12pt plus 3pt minus 12pt
  
  \ifchorded
    \beginverse*
       {\nolyrics Előjáték: \[F] \[C] \[Dm] \[A]\[A#-C] \[F] \[A] \[Dm]}
    \endverse
  \fi

  \beginchorus
    \[Dm]Jézus a mi \[A]oltalmunk, \[Dm]erőss\_{é}\[C]günk,
    Ha \[F]ránk szakad \[C]minden baj, \[Dm]mégse \[A]fé\[Dm]lünk!
    \[Dm]Ő a mi \[A]páncélunk, \[Dm]erős p\_{a}j\[C]zsunk,
    Ha \[F]ránk tör az \[C]ellenség, \[Dm]Benne \[A]bí\[Dm]zunk!
  \endchorus

  \beginverse\memorize
    \[F]Ez a föld szét\[C]hullik, \[Dm]minden hegy el\[A]omlik,
    \[A#]De egy \[F]kőszikla \[C]örökké \[F]áll!
    \[F]Tajtékzik a \[C]tenger, \[Dm]gyilkol a sok \[A]fegyver,
    \[A#]Né\[C]künk \[F]mégsem árt \[A]a ha\[Dm]lál!
  \endverse

  \beginverse*
    Refrén
  \endverse

  \beginverse
    ^Eladó már ^minden, ^örömöd még ^sincsen,
    ^Ha ez a ^forrás ^nem árad ^rád!
    ^Ő az, aki ^ingyen ^már od'adott ^mindent,
    ^Ér^tünk ^áldozva ^Önma^gát!
  \endverse

  \beginverse*
    Refrén
  \endverse

  \beginverse
    ^Fogták, elí^télték ^és megfeszí^tették
    ^Isten ^Fiát mi^helyet^tünk!
    ^Nincs már \_{í}^télet, ^vád sohasem ^érhet,
    ^Vé^re ^eltörölt ^minden ^bűnt!
  \endverse

  \beginverse*
    Refrén
  \endverse

\endsong

\beginsong{Jézus életem, erőm, békém}
[
  by={Taizé / Jacques Berthier}
]
  \beginverse*
    \[Dm]Jé\[C]zus \[F]életem, erőm, \[A#]bé\[C]kém,
    \[Dm]Jé\[C]zus \[F]társam, \[Dm]örö\[C]möm,
    Benned \[A#]bízom, \[A]Te vagy az \[Dm]Úr;
    Már \[C]nincs mit \[F]félnem, mert \[A#]bennem \[C]élsz,
    Már \[Am]nincs mit \[Dm]félnem, mert \[A#]ben\[C]nem \[F]élsz.
  \endverse

\endsong

\beginsong{Jézus néz rám}[]

  \ifchorded
    \beginverse*
      {
        \nolyrics Előjáték:
        \[C] \[Dm] \[Am] \[E7] \[Dm] \[G7] \[C] \[G7]
        \[C] \[Dm] \[Am] \[E7] \[Dm] \[G7] \[C]
      }
    \endverse
  \fi

  \beginverse
    \[C]Jézus \[Dm]néz rám a \[Am]két szemed\[E7]ből, arra \[Dm]kér, arra \[G7]vár, hogy megért\[C]sem, \[G7]
    \[C]Jézus \[Dm]néz rám a \[Am]két szemed\[E7]ből, arra \[Dm]kér, arra \[G7]vár, hogy szeres\[C]sem.
  \endverse

  \beginchorus
    /: \[C]Jézus \[H7]arra \[E]kér, mert \[C]bennünk \[H7]emberekben \[E]él,
    Hogy \[C]meglás\[C7]sam és \[F]szeressem, ha \[C]eljön hozzám \[G7]valaki\[C]ben. \[G$^{\;(csak\;ismétlésnél)}$]{\ifchorded\hspace{0.4cm}\fi:/ \rep{2}}
    % vagy inkább \quad ?
  \endchorus

\endsong

\input{songs/jezus-var-terad.tex}
\beginsong{Jó az Úrban bizakodni}
[
  by={Taizé / Jacques Berthier}
]

  \beginverse*
    \[Dm]Jó az Úrban \[A]bizakodni, {\[Dm]j}\[C]{\_{ó}} az \[F]Úr.
    \[Gm]Re\[C]mélj, és \[F]b{\_{í}}zz \[C]Ben\[Dm]ne, {\[A#]j}\[C]{\_{ó}} az \[Dm]Úr!
  \endverse

\endsong

\beginsong{Jöjj, hívunk, jöjj}[]

  \ifchorded
    \beginverse*
       {\nolyrics Előjáték: \[E]\[C#m] \[F#m] \[A] \[H] \[A] \[H]}
    \endverse
  \fi

  \beginverse\memorize
    \[E]Jöjj, hívunk, \[C#m]jöjj! Szentlé\[F#m]lek, a szí\[A]vünk\[H]be \[A]á\[H]radj!
    \[E]Jöjj, hívunk, \[C#m]jöjj! Szentlé\[F#m]lek, a szí\[A]vünk\[H]be \[A]á\[H]radj!
  \endverse

  \beginchorus
    Úgy \[A]v\_{á}\[E]runk és ki\[A]{\_{á}}l\[E]tunk:
    /: Áradj \[H]szét! \echo{Vigasztaló}
    \_{Á}radj \[A]szét! \echo{Szabadító}
    Áradj \[E]szét! :/ \rep{2}
  \endchorus

  \beginverse
    ^Jöjj, hívunk, ^jöjj! Mint fo^lyó, amely ^túl^á^ra^dó!
    ^Jöjj, hívunk, ^jöjj! Mint fo^lyó, amely ^túl^á^ra^dó!
  \endverse

  \beginverse*
    Refrén
  \endverse

\endsong

\beginsong{Jöjj, itt az idő}[]

  \ifchorded
    \beginverse*
       {\nolyrics Előjáték: \[D] \[D4] \[D] \[D4]}
    \endverse
  \fi

  \beginverse\memorize
    \[D]Jöjj, itt az idő, hogy \[D4]éb\[D]redj!
    \[A]Jöjj, itt az idő, hogy \[Em]áld\[F#m]juk \[G]Őt!
    \[D]Jöjj, szívedet add az \[D4]Úr\[D]nak!
    \[A]Jöjj, imádd az Istent \[Em]úgy, \[F#m]ahogy \[G]vagy!
  \endverse

  \beginverse*
    Hát \[D]jö\[D2 D4 D]jj!
  \endverse

  \beginchorus
    /: \[G]Minden nyelv megv\_{a}llja, hogy \[D]Ő az Úr! \[G]És mindenki t\_{é}rdet \[D]hajt!
    De a \[G]legnagyobb kincs mégis \[Hm]azokra vár, kik \[Em]most Őt válasz\[A4]tják! \[A] :/ \rep{2}
  \endchorus

  \beginverse*
    1. versszak
  \endverse

  \beginverse*
    Hát \[D]jöjj! \[G] Ó \[D]jöjj! \[G] Hát \[D]jöjj!
  \endverse

\endsong

\beginsong{Kész az én szívem}[]

  \ifchorded
    \beginverse*
       {\nolyrics Előjáték: \[Dm] \[Dm]}
    \endverse
  \fi

  \beginverse\memorize
    Kész az \[Dm]én sz\_{í}vem, Uram, \[Gm]Néked zeng a \[Dm]dalom.
    Elő hát hárfa, lant, hadd kel\[C]tsem fel a h\_{a}j\[A]nalt! \[A4] \[A]
  \endverse

  \beginchorus
    /: Szent az \[Dm]Úr, jó az Úr, hirde\[Gm]tem a népek \[Dm]között,
    Mert az Ő kegyel\[Gm]me felhőkig \[Dm]ér. \[A-A4$^{\;(csak\;ismétlésnél)}$]{\ifchorded\hspace{1cm}\fi:/ \rep{2}}
  \endchorus

  \beginverse
    Igaz^ság, hatalom igéd ^minden egyes ^szava.
    Szívemben rejtem el, soha ^el ne fel\_{e}j^tsem! ^ ^
  \endverse

  \beginchorus
    /: Szent az \[Dm]Úr, jó az Úr, hirde\[Gm]tem a népek \[Dm]között,
    Mert az Ő hűsé\[Gm]ge az égig \[Dm]ér. \[A-A4$^{\;(csak\;ismétlésnél)}$]{\ifchorded\hspace{1cm}\fi:/ \rep{2}}
  \endchorus

\endsong

\input{songs/ket-kezemben.tex}
\input{songs/kicsiny-kis-fenyemmel.tex}
\beginsong{Király vagy}[]

  \beginverse*
    /: \[D]Király vagy, \echo{Király vagy}
    \[A]Király v\_{a}gy, \echo{Király v\[Hm]{\_{a}}gy}
    Jézus Ki\[G]rály. :/ \rep{2}
  \endverse

  \beginverse*
    /: Most \[D]felemeljük a szívünket,
    Most \[A]felemeljük kezeinket,
    \[Hm]Trónod elé járul\_{u}nk, imádva \[G]légy! :/ \rep{2}
  \endverse

\endsong


\beginsong{Köszönöm, Jézus}[]

  \beginverse\memorize
    \[D] Köszönöm, J\_{é}zus, köszönöm, J\_{é}zus,
    Köszönöm, U\[G]ram, hogy szeretsz en\[D]gem!
    Köszönöm, \[F#m]J{\_{é}}zus, köszönöm, \[Hm]J{\_{é}}\[G]zus,
    Köszönöm, U\[D]ram, \[A] hogy szeretsz en\[D]gem! \[G] \[D]
  \endverse

  \beginverse
    ^ A Golgotára ment, és meghalt \_{é}rtem;
    Köszönöm, U^ram, hogy szeretsz en^gem!
    A Golgo^tára ment, és meghalt ^{\_{é}}r^tem;
    Köszönöm, U^ram, ^ hogy szeretsz en^gem! ^ ^
  \endverse

  \beginverse
    ^ A harmadik napon Ő f\_{e}ltámadott;
    Köszönöm, U^ram, hogy szeretsz en^gem!
    A harma^dik napon Ő f\_{e}l^táma^dott;
    Köszönöm, U^ram, ^ hogy szeretsz en^gem! ^ ^
  \endverse

  \beginverse
    ^ Ó, allel\_{u}ja, ó, allel\_{u}ja!
    Köszönöm, U^ram, hogy szeretsz en^gem!
    Ó, alle^l\_{u}ja, ó, alle^l{\_{u}}^ja!
    Köszönöm, U^ram, ^ hogy szeretsz en^gem! ^ ^
  \endverse

\endsong

\beginsong{Krisztus él bennem}[]

  \ifchorded
    \beginverse*
       {\nolyrics Előjáték: \[G] \[D] \[C] \[C-D]}
    \endverse
  \fi

  \beginchorus
    \[G]Krisztus él bennem, \[D]szívem virágzik.
    A \[Em]sötét nem győz, hát \[C]nem kell pánik.
    \[G]Jézus \[D]életét \[C]adta, hogy élj! \[D]

    \vspace{0.3cm}

    \[G]Krisztus él benned, \[D]ne légy lekvár,
    Hisz \[Em]oly sok kincs van, mi \[C]téged is vár!
    Hát \[G]kelj fel, \[D]indulj és \[C]Krisztussal járj!

    \vspace{0.3cm}

    /: \[D]Na-na-na, \[G]na na na \[D]na, na-na-na, \[Em]na na na \[C]na,
    na-na-na, \[G]na na na, \[D]na na na, \[C]na na! :/ \rep{2} \[D$^{\;(a\;legvégén\;nincs)}$]
  \endchorus

  \beginverse\memorize
    Az \[Em]Örömhír fénye \[C]ragyogjon éle\[D]tedben!
    A \[G]szeretet tüze \[Am]égjen a szíved\[D]ben!
    Hogy \[Hm]lássa meg a vi\[Em]lág, te is \[C]mondd velünk, \[Am]éld meg és add to\[D]vább! \[D4] \[D]
  \endverse

  \beginverse*
    Refrén
  \endverse

  \beginverse
    ^Hiába háborog ^alattunk a ten^ger,
    ^Jézus szemébe ^nézve félnünk ^nem kell.
    Ha ^Ő bennünk ^él, előttünk ^világos, ^ tiszta a ^cél. ^ ^
  \endverse

  \beginverse*
    Refrén
  \endverse

\endsong

\beginsong{Látod, újra este van}[]

  \beginverse\memorize
    \[C] Látod, újra este van, Hozzád száll most \[C/H]halk sza\[Am]vam,
    \[Dm] Mindenségnek Iste\[F]n\[G]{\_{e}}, Téged dicsér éne\[C]kem.
  \endverse

  \beginverse
    ^ Köszönöm ezt a napot, annyi szépet s ^jót ho^zott.
    ^ Köszönöm a napsü^t\_{é}^st, fényét szórta szerte^szét.
  \endverse

  \beginchorus
    \[C]Kérlek, ó jó Uram, \[Dm]add meg, fényeddé \[F]váljak, jeled le\[C]hessek!
    \[C]Kérlek, Szentl\_{e}lked \[Dm]küldd el, \_{Ő} ve\[F]zessen \[G]végig élete\[C]men!
  \endchorus

  \beginverse
    ^ Uram, mindent köszönök, bánatot é^s örö^möt,
    ^ Legyen rajtam szent ke^z\_{e}^d, óvja, védje éle^tem!
  \endverse

  \beginverse
    ^ Szeretném, ha szeretnél, minden szavad ^érte^ném,
    ^ Életemben követ^n{\_{é}}^m, s végül Hozzád érkez^ném.
  \endverse

  \beginverse*
    Refrén

\beginsong{Lelkem, áldd az Urat}[]

  \versesep=12pt plus 3pt minus 12pt

  \ifchorded
    \beginverse*
       {\nolyrics Előjáték: \[D] \[D4] \[D]}
    \endverse
  \fi

  \beginchorus
    Lelkem, \[G]áldd az Ur\[D]at, \[A/C#]{\_{á}}ldd az U\[Hm]rat,
    \[G]Egyedül \[D]Őt im\[A4]ádd! \[A]
    Zengj úgy, \[G]mint soha \[Hm]még, \[G]Őt \[A]di\[Hm]csérd!
    A l\[G]elkem csak T\[A]éged \[G/D]áld! \[D] \[G/D] \[D]
  \endchorus

  \beginverse\memorize
    Egy \[G]új nap \[D]kél, és már \[A]ébred \[Hm]minden,
    \[G] Itt az id\[D]ő, hogy \[A]áldja\[Hm]lak,
    És \[G]bármi \[D]ér, bármit r\[A]ejt az út el\_{ő}t\[Hm]tem,
    \[G] Hadd zengjek \[D]este is majd \[A4]há\[A]lá\[D]val. \[D4] \[D]
  \endverse

  \beginverse*
    Refrén
  \endverse

  \beginverse
    ^ Kegyelmed ^nagy, és tü^relmed ^hosszú,
    ^ Szerelmed ^mély, és az ^égig ^ér,
    Nem ^szűnök ^énekelni ^jósá^g{\_{o}}dról,
    ^ Tízezer ^okom van, hogy ^zeng^jek ^még. ^ ^
  \endverse

  \beginverse*
    Refrén
  \endverse

  \beginverse
    ^ És egy nap ^majd ez a ^test el^gyengül,
    ^ Erőm el^hagy, s utam ^véget ^ér,
    A ^lelkem ^mégis Téged ^áld vég ^n{\_{é}}lkül:
    ^ Tízezer ^évig és mind^örö^kké! ^ ^ ^
  \endverse

  \beginverse*
    Refrén \rep{2}
  \endverse

\endsong

\input{songs/lelkem-az-urat-dicserd.tex}
\beginsong{Loyolai Szent Ignác imája}
[
  index={Fogadd el}
]

  \ifchorded
    \beginverse*
       {\nolyrics Előjáték: \[C] \[G] \[Am] \[F] \[C] \[G] \[C] \[G7]}
    \endverse
  \fi

  \beginverse\memorize
    \[C]Fogadd el, \[G]Uram, \[Am]szabadságo\[F]mat, \[C]fogadd el \[G]eg{\_{é}}\[C]szen! \[G]
    \[C] Vedd értel\[G]memet, \[Am]akarato\[F]mat, és \[C]emlé\[G]kez{\_{é}}\[C]sem!
  \endverse

  \beginverse
    ^Mindazt, amim ^van ^és ami ^vagyok, ^Te adtad ^{\_{i}}n^gyen.^
    ^Visszaadok, ^Uram, ^visszaa^dok ^egyszerre ^m\_{i}n^dent.
  \endverse

  \beginchorus
    \[Am]Legyen \[E7]fölöttük \[Am]korlátlan \[F]úr \[C]rendel\[G]kezé\[C]sed!
    Csak \[Am]egyet \[E7]h{\_{a}}gyj meg \[Am]ajándéko\[D7]dul: \[G]szeretnem \[F]T{\_{é}}\[Em]ged! \[G7]
  \endchorus

  \beginverse
    ^Csak a ^szeretet ^maradjon e^nyém a ^kegye^l{\_{e}}m^mel, ^
    S ^minden, de ^minden ^gazdagság e^nyém, más ^semmi ^n{\_{e}}m ^kell!
  \endverse

  \ifchorded
    \beginverse*
       {\nolyrics Utójáték: \[C] \[G] \[Am] \[F] \[C] \[G] \[C]}
    \endverse
  \fi

\endsong

\beginsong{A mélyből Hozzád száll szavam}[]

  \versesep=12pt plus 3pt minus 9pt

  \ifchorded
    \beginverse*
      {\nolyrics Előjáték: \[H] \[A] \[G] \[F#]}
    \endverse
  \fi

  \beginverse\memorize
    A \[Em]mélyből \[H7]Hozzád \[Em]száll sza\[Am]vam, \[G]Uram, \[H7]irgal\[Em]mazz!
    A \[Em]bajban \[H7]lelkem \[Em]társta\[Am]lan, \[G]Uram, \[H7]irgal\[Em]mazz!
  \endverse

  \beginchorus
    \[Em!]Segíts! \echo{Segíts!} \[Am!]Ne hagyj! \echo{Ne hagyj!}
    \[Em]Nézd, \[H7]ránk tör \[Em]a bűn m\_{á}r!
    \[Em!]Te adj \echo{Te adj} \[Am!]erőt! \echo{erőt!}
    \[G]Uram, \[H7]irgal\[Em]mazz!
  \endchorus

  \beginverse
    Ha ^ajkam ^vétett ^elle^ned, ^Krisztus, ^kegyel^mezz!
    A ^jóhírt ^zengjem ^csak Ve^led, ^Krisztus, ^kegyel^mezz!
  \endverse

  \beginchorus
    \[Em!]Segíts! \echo{Segíts!} \[Am!]Ne hagyj! \echo{Ne hagyj!}
    \[Em]Nézd, \[H7]ránk tör \[Em]a bűn m\_{á}r!
    \[Em!]Te adj \echo{Te adj} \[Am!]erőt! \echo{erőt!}
    \[G]Krisztus, \[H7]kegyel\[Em]mezz!
  \endchorus

  \beginverse
    Ha ^arcod ^fényét ^elfe^dem, ^Uram, ^irgal^mazz!
    Te ^újból ^hívjál, ^Iste^nem, ^Uram, ^irgal^mazz!
  \endverse

  \beginchorus
    \[Em!]Segíts! \echo{Segíts!} \[Am!]Ne hagyj! \echo{Ne hagyj!}
    \[Em]Nézd, \[H7]ránk tör \[Em]a bűn m\_{á}r!
    \[Em!]Te adj \echo{Te adj} \[Am!]erőt! \echo{erőt!}
    \[G]Uram, \[H7]irgal\[Em]mazz!
  \endchorus

\endsong

\input{songs/mennek-az-asszonyok.tex}
\beginsong{Mert tudom}[]

  \beginverse*\memorize
    Mert tu\[D]d{\_{o}}m, semmi \[G]nem szakaszthat \[D]el az én \[G]Uram szerelmé\[D]től.
    S\_{e}m \_{é}let, sem ha\[G]lál nem választ \[D]el, sem \[G]m{\_{a}}, s\_{e}m \[D]holnap.
  \endverse

  \beginchorus
    Mert az \[D]angyalok \[G]visznek Őhoz\[D]zá,
    S a po\[F#m]kol e\[Hm]r{\_{ő}}i sem \[G]állítnak \[A]meg. \[A4] \[A]
    Nem \[G]számít, ha hegyte\[D]tőn állok meg, vagy a \[Hm]sík ten\[G]ger part\[A]ján,
    Mert az \[D]{\_{Ú}}r alkotta \[Hm]m{\_{i}}ndezt, és \[G]minden az Ö\[A]vé.
  \endchorus

  \beginverse*
    És tu^d{\_{o}}m, semmi ^nem szakaszthat ^el az én ^Uram szerelmé^től.
  \endverse

\endsong

\beginsong{Minden, mi él}[]
  \versesep=12pt plus 3pt minus 8pt

  \beginchorus
    \[D]Minden, mi él, csak Téged h\_{i}rdet,
    \[Hm]Minden dicsér, mert mind a m\_{ű}ved,
    \[G]Azzal, hogy él, ezt zengi N\_{é}ked:
    \[A]Dicsérlek én, \[A7]dicsérlek T\_{é}g\_{e}d!
  \endchorus

  \beginverse\memorize
    \[D]Dicsér az ég, Nap, Hold és csillagok,
    \[Hm]Fény és sötét, nap, éj és hajnalok.
    \[G]Dicsér a szél, felhő és hóvihar,
    \[A]A víz s a tűz, \[A7]megannyi tiszta dallal!
  \endverse

  \beginverse*
    Refrén
  \endverse

  \beginverse
    ^Dicsér a Föld, dicséri szent neved,
    ^Mint jó anyánk, táplál s ad eledelt.
    ^Virág, gyümölcs, zöld fű, fa, hegyvidék,
    ^Tó és folyó, ^síkság és büszke b\_{é}rc!
  \endverse

  \beginverse*
    Refrén
  \endverse

  \beginverse
    ^A nagy világ létével énekel,
    ^Szavunkra vár, hogy hangja dal legyen.
    ^Zengjük tehát ég és föld énekét,
    ^Zengjük velük: ^nagy Isten áldott l\_{é}gy!
  \endverse

  \beginverse*
    Refrén \echo{\[D]Dicsérlek én!}
  \endverse

\endsong

\beginsong{Minden nap}
[
  by={Pintér Béla},
  index={Mit is mondhatnék}
]

  \ifchorded
    \beginverse*
      {\nolyrics Előjáték: \[D] \[A] \[F#m] \[E]}
    \endverse
  \fi

  \beginverse\memorize
    \[D] Mit is \[A]mondhatnék, \[F#m] Te vagy, ki \[E]megváltottál,
    \[D] Drága \[A]véreddel \[F#m] hófe\[E]hérre mostál,
    \[D] Minden\[A]em lettél, \[F#m] Tiéd az \[E]egész lényem,
    \[D] Ég a \[A]tűz bennem, \[F#m] csak \[E]Érted élek \[A]m{\_{á}}r! \[D] \[F#m] \[E]
  \endverse

  \ifchorded
    \beginverse*
      {\nolyrics Közjáték: \[A] \[D] \[F#m] \[E]}
    \endverse
  \fi

  \beginverse
    ^ Minden ^nap, Uram, ^ ráál^lok igédre,
    ^ Téged ^hívl\_{a}k, ^ légy a ^segítségem,
    ^ Vezess ^át engem, ^ tovább a ^keskeny úton,
    ^ Add, hogy ^életem ^ a világban ^tüzet gy\_{ú}jt^{s\_{o}}n! ^ ^ ^
  \endverse

  \beginchorus
    \[A] Minden \[D]nap csak \[F#m]Érted \[E]élek.
    \[A] Minden \[D]nap kö\[F#m]vetlek \[E]Téged.
    \[A] Minden \[D]nap Ve\[F#m]led já\[E]rok, U\[A]ram. \[D] \[F#m] \[E]
  \endchorus

  \beginverse*
    /: \[A]Érted \[D]élek \[F#m]minden \[E]n{\_{a}}p! :/ \rep{3}
    U\[A]ram! \[D] \[F#m] \[E]
  \endverse

  \beginverse*
    Refrén \rep{2} \[A]
  \endverse

\endsong

\beginsong{Mint az asszony}[]

  \beginverse*\memorize
    Mint az \[A]asszony az Ő \[E/G#]ruhája szegé\[F#m]lyét, \[F#m/E]
    Megragad\[D]juk \[D/C#] az Úr je\[Hm]lenl{\_{é}}\[E]tét.
    Mint a \[A]vak koldus az \[E/G#]{\_{ú}}t sz\_{é}\[F#m]lén, \[F#m/E]
    Kiál\[D]tunk, \[D/C#] s fu\[Hm]tunk az Úr e\[E]lé.
  \endverse

  \beginchorus
    És \[A]hirtelen \[E] egy mennyei \[D]kéz megáld, \[D/C#]
    \[Hm]Jézus az, \[E] aki megse\[A]gít! \[E]
    És \[A]hirtelen \[E] a mennyei \[D]erő leszáll, \[D/C#]
    \[Hm]Jézus jön \[E] s meggy\_{ó}\[A]gyít!
  \endchorus

\endsong

\beginsong{Most nem sietek}
[
  by={Túrmezei Erzsébet}
]

  \beginverse\memorize
    \[Em]Most nem sie\[G]tek, most nem roha\[D]nok,
    \[Em]Most nem terve\[G]zek, \[C]most nem aka\[D]rok,
  \endverse

  \beginchorus
    \[G]Most nem teszek \[D]semmit sem,
    Csak \[C]engedem, hogy \[D]szeressen az \[G]Ist\_{e}n.
    \[G]Most nem teszek \[D]semmit sem,
    \[G]Csak engedem, hogy \[H7]szeressen az \[Em]Isten.
  \endchorus

  \beginverse
    ^Most m\_{e}gnyug^szom, most elpihe^nek,
    ^Békén, szaba^don, ^mint gyenge gye^rek,
    Nem, nem teszek ...
  \endverse

  \beginverse
    ^S míg ölel a ^fény, és ölel a ^csend,
    ^És árad be^lém, ^és újjáte^remt,
    Míg nem teszek ...
  \endverse

  \beginverse
    ^Új gyümölcs te^rem, másoknak te^rem,
    ^Érik csende^sen ^erő, győze^lem,
    Ha nem teszek ...
  \endverse

\endsong

\input{songs/mustarmag.tex}
\beginsong{Napfivér, Holdnővér}[]

  \ifchorded
    \beginverse*
      {\nolyrics Előjáték: \[D4] \[D] \[D4] \[D]}
    \endverse
  \fi

  \beginverse\memorize
    \[D]Fivérem \[F#m]Nap és \[G]n{\_{ő}}\[A]vérem \[D]Hold,
    \[Hm]Oly' ritkán \[F#m]látlak, s \[G]hallom \[A]hango\[D]tok.
    \[Hm] Nyomaszt a \[F#m]sok \[G]gyötre\[A]lem és \[D]gond.
  \endverse

  \ifchorded
    \beginverse*
      {\nolyrics Közjáték: \[Hm] \[F#m] \[G] \[A] \[D]}
    \endverse
  \fi

  \beginverse
    ^Fivérem ^Szél és ^Le^veg\_{ő}^ég,
    ^ Nyisd ki sze^mem, hogy ^lássam, ^ami ^szép!
    ^ Körülö^lel a ^ragyo^gás, dics^fény.
  \endverse

  \beginchorus
    \[Hm]Mert Isten \[F#m]m{\_{ű}}ve \[G]minden \[A]teremt\[D]mény,
    \[Hm] Érzem jó\[F#m]ságát, és \[G]szívem újra {\[A4 A]é}l!
  \endchorus

  \beginverse
    ^Fivérem ^Nap és ^n{\_{ő}}^vérem ^Hold,
    ^Most végre ^látok, s ^hallom ^hango^tok!
    ^ Megölel^ném az ^egész ^vil\_{á}^got!
  \endverse

  \beginverse*
    Refrén
  \endverse

  \beginverse*
    ^Fivérem ^Nap és ^n{\_{ő}}^vérem ^Hold,
    ^Most végre ^látok, s ^hallom ^hango^tok!
    ^ Megölel^ném az ^egész ^vil\_{á}\[Gm]got! \[Gm7] ^
  \endverse

\endsong

\beginsong{Ne félj, mert megváltottalak}[]

  \beginchorus
    /: \[D]Ne félj, mert \[G]megváltotta\[A]lak,
    Neveden \[D]szólította\[Hm]lak,
    Karja\[Em]imba zárta\[A]lak,
    Örökre \[A7]enyém \[D]vagy! :/ \rep{2}
  \endchorus

  \beginverse\memorize
    \[A]Viruló réteken \[D]át \[G]hűs forrás fe\[A]lé vezetlek,
    \[Hm]Pásztorod va\[Em]gyok, elveszni \[A]senkit nem ha\[D]gyok.
    Karom fe\[Hm]léd t\_{á}\[Em]rom, kiárad \[G]áld\_{á}\[D]som.
    \[Hm]Nem rejtőzöm \[Em]el, szívem a \[A]szívednek fe\[D]lel,
    Amikor \[Hm]úgy \_{é}r\[Em]zed, nyomaszt a \[A7]{\_{é}}\[D]let.
  \endverse

  \beginverse*
    Refrén
  \endverse

  \beginverse
    ^Nem taszítalak ^el, ^amikor vétke^zel,
    Irgalmat ^lelsz a szívem^ben, örök fe^léd a hűsé^gem,
    Amerre ^jársz, v\_{é}^dlek, nyomodba ^l{\_{é}}^pek.
    ^Nem rejtőzöm ^el, szeretet^lángom átö^lel,
    Ne félj, ha ^éjb{\_{e}}n ^jársz, hidd, hogy a ^fény v\_{á}r ^rád!
  \endverse

  \beginverse*
    Refrén
  \endverse

\endsong

\beginsong{Ne félj, ne aggódj}
[
  by={Taizé / Jacques Berthier}
]

  \beginverse*
    \[Am]Ne félj, ne \[Dm7]aggódj, \[G]ne sírj, ne \[Cmaj7]bánkódj:
    \[F]Ha tiéd \[Dm6]Isten, \[E]tiéd már \[Am]minden.
    \[Am]Ne félj, ne \[Dm7]aggódj, \[G]ne sírj, ne \[Cmaj7]bánkódj:
    \[F]Elég \[Dm6]Ő \[E]né\[Am]ked.
  \endverse

\endsong

\beginsong{Néked hódolok}[]

  \beginverse*
    \[A]Néked hódo\[Hm]lok Jé\[E]zus, dicső Ki\[A]rály! \[D] \[E]
    \[A]Néked hódo\[Hm]lok Jé\[E]zus, dicső Ki\[A]rály!
    Ujjongj, da\[Hm]lolj szív\[E]{\_{e}}m, és imádjad \[A]Őt!
    \[F#m] Ujjongj, da\[Hm]lolj szív\[E]{\_{e}}m, és imádjad \[A]Őt!
  \endverse

\endsong

\input{songs/nezz-testver-fel.tex}
\input{songs/nincs-mas-isten.tex}
\beginsong{Oly' jó áldani}[]

  \beginverse*
    \[F]Szívemből \[C]ünnepi \[F]ének á\[C]rad,
    \[G]Háladalt \[G7]éneklek az \[F]ég Urá\[C]nak!
    \[G]Ujjongva \[G7]tör fel a \[F]hála ben\[C]nem,
    \[Dm]Jósága körülvesz \[G]szárnyként \[G7]en\[C]gem. \[Dm] \[G] \[C]
  \endverse

  \beginverse*
    \[G]Adjatok \[G7]hálát az \[C]Istennek!
    \[G7]Adjatok hálát az \[C]Istennek!
    \[C7]Ő a mi teremtő \[F]Kirá\[G]lyunk,
    \[D]Alle-allelu\[G]ja, \[F]allelu\[C]ja, \[F]allelu\[C]ja! \[H7]
  \endverse

  \beginverse*
    /: \[E]Hála legyen, hála legyen, hála! \[A]Áldott az \[E]Úr! :/ \rep{3}
  \endverse

  \ifchorded
    \beginverse*
      {\nolyrics Átvezető: \[D] \[A] /: \[A] \[E] \[A] \[E] :/ \rep{2}}
    \endverse
  \fi

  \beginverse\memorize
    \[A]Oly' jó ál\[E]dan\[A]i az \[E]Ur\[A]at, \[A] \[E] \[A] \[E]
    \[E]Ujjongó\[A]an \[E]jó Is\[A]ten\[E]nek zen\[A]ge\[D]ni zsoltáro\[A]kat. \[A] \[E] \[A] \[E]
  \endverse

  \beginverse
    ^Már kora reg^ge^len jó^sá^gát, ^ ^ ^ ^
    És ^éjszaká^kon ^át hir^de^tem szün^te^len nagy irgal^mát. ^ ^ ^ ^
  \endverse

  \beginverse
    ^Pengő hang^sze^ren, lan^to^kon, ^ ^ ^ ^
    ^Ujjongó^an ^jó Is^ten^nek zen^ge^ni hála^szót. ^ ^ ^ ^
  \endverse

  \beginverse
    ^Az öröm hang^ja^i szól^ja^nak, ^ ^ ^ ^
    ^Gitár húr^ja^in magasz^ta^lom uj^jong^va csodái^dat! ^ ^ ^ ^
  \endverse

  % \ifchorded
  %   \chordsoff
  % \fi

    \beginverse
      ^Uram, jó voltál énhozzám,
      Életemben nagy dolgokat műveltél sok év során.
    \endverse

    \beginverse
      Sziklám, oltalmam Te voltál,
      Oly' irgalmasan jó, gondoskodó Istenem, égi Atyám.
    \endverse

  % \ifchorded
  %   \chordson
  % \fi

  \beginverse
    ^Hála, di^cső^ség Ne^ved^nek, ^ ^ ^ ^
    ^Most és ö^rök^ké az ^é^gen, a ^Föl^dön! \_{A}^men. ^ ^ ^ ^ \[A]
  \endverse

\endsong

\beginsong{Ott lélegzel a fákban}[]

  \ifchorded
    \beginverse*
       {\nolyrics Előjáték: \[D] \[A] \[A7] \[D]}
    \endverse
  \fi

  \beginverse\memorize
    \[D]Ott lélegzel a \[A]f{\_{á}}k\[D]ban, \[G]ott mosolyogsz a vi\[A]r{\_{á}}gban,
    \[A]Ott vagy a szélben, a \[D]fr{\_{i}}ssben, ó, boldog, b\_{o}ldog \[A7]{\_{I}}s\[D]ten!
  \endverse

  \beginverse
    ^Szép sorsomat Te ^sz{\_{ő}}^tted, ^jó leborulni e^l{\_{ő}}tted,
    ^Te vagy az útnak a ^v{\_{é}}ge, sz\_{e}nt mélység, \_{á}ldott ^b{\_{é}}^ke!
  \endverse

\endsong

\input{songs/o-az-ur.tex}
\beginsong{Péter, ne sírj}[]

  \beginverse\memorize
    \[Em]Sírásodat hagyd nyugodtan abba,
    A kapuk mögött vissza\[Em/F#]állt a \[G]rend!
    \[Am]Ne sírj, oly' lágyan jár a \[Em]sz{\_{é}}l
    \[C]Róma s a \[D]Via Appia fe\[Em]lett!
  \endverse

  \beginchorus
    \[Am]Péter, ne \[Em]sírj, mert \[H7]Jézus megbo\[Em]csát,
    \[Am]Péter, ne \[Em]s{\_{í}}rj, hisz' \[H7]tudod, hogy megbo\[Em]csát!
  \endchorus

  \beginverse
    ^Gyengeségünk néha úgy esik ránk,
    Mint kisgyermekre esik a ^féle^lem.
    ^Gyengeségért mindig bünte^t{\_{é}}s jár,
    De ^büntetése ^nem sújt, fele^mel!
  \endverse

  \beginverse*
    Refrén
  \endverse

  \beginverse
    ^Reménytelen a bűn után az élet,
    Hogy mégis élsz, már alig ^hiszed ^el,
    ^Alig hiszed, hogy nem lett akkor ^v{\_{é}}ge,
    Csak ^egyet tudsz, a ^hulló könnye^ket!
  \endverse

  \beginverse*
    Refrén \rep{2}
  \endverse

\endsong

\input{songs/rajta-dicserjetek-az-ur-nevet.tex}
\beginsong{Rejts most el}
[
  by={Reuben Morgan, fordította Labadics Kriszta}
]

  \beginverse\memorize
    \[C]Rejts \[G/H]most \[Am]el a \[F]szár\[D/F#]nyad a\[G]lá,
    E\[C]rős \[G]kéz\[F]zel \[Dm]takarj \[D]be \_{e}n\[G]gem!
  \endverse

  \beginchorus
    \[C]Tenger tombol, \[F]zúg, süv\[G]{\_{í}}t a sz\[C]él,
    \[C]Te emelsz fel \[F]a vih\[G]ar fö\[Am]lé!
    \[C]Uralkodsz hul\[F]lámok h\[G]{\_{a}}bja\[C]in,
    \[C]Szívem nem f\[F]él, Benn\[G]ed rem\[Am]él!
    \[C]Szívem nem f\[F]él, Benn\[G]ed rem\[F]él! \[G]
  \endchorus

  \beginverse
    Csak ^Is^ten^ben bízz, ^én ^l{\_{e}}l^kem!
    ^Mert ^Ő ^él, és ^nagyobb ^mind\_{e}n^nél!
  \endverse

  \beginverse*
    Refrén  \[C]
  \endverse

\endsong

\input{songs/sotet-fenn-az-eg.tex}
\beginsong{Számtalan a csoda}[]

  \beginverse*\memorize
   \[C]Számtal\[G]an a csoda, \[Dm] amit m\[Am]űvelt az Úr!
    \[C]Nézd, mennyi \[F]virág nyílik a \[C]sziklák közt a \[G]földön!
    Az \[C]Ő sza\[G]va az élet, \[Dm] minden \[Am]igéje tűz,
    \[C]Bennünk akar \[F]továbbélni \[C]mindenütt a \[G]földön,
    Mert \[Am]Ő az a bővizű, \[Dm]életadó, tiszta \[Em]forrás:
    Ki ebből mer\[F]ít, soha nem szomjazik \[G]már!
  \endverse

 \beginverse*
    ^Úgy, mint a ^mag, mely ^földbe hullik ^és kihajt:
    ^úgy termi ^élő Igéd ^bennünk égi ^gyümölcsét.
    ^Fény és re^mény, ^béke, öröm, ^szabadság:
    ^Jézussal ^köztünk vagyunk
    Az ^egész Földön ^egy család!
  \endverse

  \beginverse*
    Számtalan a csoda...
  \endverse

\endsong

\beginsong{Szentlélek, jöjj}[]

  \versesep=12pt plus 3pt minus 12pt

  \ifchorded
    \beginverse*
       {\nolyrics Előjáték: \[Em] \[D] \[Hm] \[Em] \[D] \[Hm] \[Em]}
    \endverse
  \fi

  \beginchorus\memorize
    \[Em]Szentlélek, jöjj, lobogó \[D]Láng!
    Szentlélek, jöjj, \[Hm]a világ \[Em]vár!
    Szentlélek, jöjj, viharos \[D]szél!
    Jöjj, \[Hm]áradj \[Em]szét!
  \endchorus

  \beginverse
    ^Jöjj el, Élő vízfor^rás,
    Jöjj, a szívünk ^Téged ^vár!
    Jöjj, ki fényt adsz lelkünk^nek,
    Jöjj, úgy ^várunk ^Rád!
  \endverse

  \beginverse*
    Refrén
  \endverse

  \beginverse
    ^Jöjj, igazság forrá^sa,
    Jöjj, imádunk ^mindnyá^jan!
    Jöjj, reményünk éleszd ^fel,
    Jöjj ke^gyelmed^del!
  \endverse

  \beginverse*
    Refrén
  \endverse

  \beginverse
    ^Jöjj, a néped gyűjtsd egy^be!
    Jöjj, az alvót ^ébreszd ^fel!
    Jöjj, a bűntől tisztíts ^meg,
    Báto^ríts min^ket!
  \endverse

  \beginverse*
    Refrén \rep{2}
  \endverse

\endsong

\input{songs/szeretem-orokke.tex}
\input{songs/tancolj-az-urnak.tex}
\input{songs/terad-var-egy-szep-orszag.tex}
\beginsong{Teremts bennem}
[
  li={Kék könyv / 208.}
]

  \beginverse*
    /: \[G]Teremts bennem \[D]tiszt\_{a} szí\[C]vet, \_{ó} U\[G]ram!
    Az \[Em]erős lelket \[D]újítsd meg ben\[G-C]nem! \[G] :/ \rep{2}
  \endverse

  \beginverse*
    /: \[Am]Ne vess el en\[D]gem a Te or\[G]cád e\[E]lől!
    \[Am]Szentlelked \[D]ne vondd meg tő\[G]lem! \[G7]
    \[Am]Támogass az \[D]engedelmes\[G]ség lelké\[Em]vel!
    \[Am7]Szabadításod \[D7]örömét add ne\[G]kem! \[G7] :/ \rep{2}
  \endverse

  \beginverse*
    \[Am7]Szabadításod \[D7]örömét add ne\[G-C]kem! \[G]
  \endverse

\endsong

\input{songs/tied-a-dicsoseg.tex}
\input{songs/tied-vagyunk.tex}
\beginsong{Uram, Tehozzád futok}[]

  \beginverse*\memorize
    Uram, \[Em]Tehozzád fu\[D]tok, élő \[Hm7]vízre szomja\[Em]zom,
    Közel\[C]séged, ó mily \[D]jó énné\[Em]kem. \[Em!]
  \endverse

  \beginverse*
    Kérlek, ^ne menj el tő^lem, légy min^dig segítsé^gem,
    Úgy kí^vánlak Té^ged, Iste^nem. ^
  \endverse

  \beginverse*
    Én pedig \[G]szüntelen remél\[D]ek, egyre \[Hm7]jobban dicsér\[Em]lek,
    Ajkam \[Am7]beszéli a Te \[Am/F#]igazságo\[H4]dat. \[H]
    Hadd le\[G]gyen most a da\[D]lom jó i\[Hm]llat oltáro\[Em]don,
    Nagyon \[C]szeretlek Té\[D]ged, Jézu\[Em]som. \[Em!]
  \endverse

  \beginverse*
    \textnote{A legvégén:}
    Nagyon \[C]szeretlek Té\[D]ged,
    Nagyon \[C]szeretlek Té\[D]ged,
    Nagyon \[C]szeretlek Té\[D]ged, Jézu\[Em]som. \[Em!]
  \endverse

\endsong

\beginsong{Az Úrra vár a szivünk}
[
  by={Taizé / Jacques Berthier}
]

  \beginverse*
    Az \[Dm]Úrra \[C]vár a szi\[F]vünk,
    Ő\[Gm]benne \[Dm]minden \[A]örö\[Dm]münk.
  \endverse

  \ifchorded
    \beginverse*
      {\nolyrics ( Köztes: \[Dm] \[A#] \[C]
      \[F] \[Dm] \[Em] \[A] )}
    \endverse
  \fi

\endsong

\beginsong{Utad vár rád}
[
  index={Mélyen a lelkedben}
]

  \ifchorded
    \beginverse*
      {\nolyrics Előjáték: \[Dm-Dm4-Dm] \[F-F2-F] \[C-C4-C] \[Dm!-A!]}
    \endverse
  \fi

  \beginverse\memorize
    \[Dm]Mélyen a \[A]lelkedben egy \[Dm]tiszta fo\[C]lyó,
    Amit \[F]áradni \[Gm]érzel, Ő a \[C]Mindenha\[Dm]tó.
    \[Dm]Isteni \[A]üzenet az \[Dm]éjben a \[C]f{\_{é}}ny,
    \[F]Királysága \[Gm]már a \[C]szívedben \[Dm]él!
  \endverse

  \beginchorus
    Utad \[F]vár \[C]rád, \[Gm]indulj \[Dm]hát!
    \[F]Nálad van a \[Gm]kincs, mire \[C]mindenki \[Dm]vágy,
    Utad \[F]vár \[C]rád, \[Gm]indulj \[Dm]hát!
    \[F]Benned az \[G]isteni \[Gm]tűz a földre \[Dm]sz{\_{á}}ll.
  \endchorus

  \beginverse
    ^Reménységed ^forrása ^égi A^ty{\_{á}}d,
    ^Szeretetének ^nincs, ki ^{\_{e}}llen^áll.
    ^Te lehetsz a ^fénylő jel a ^hegy tete^j{\_{é}}n,
    ^Te lehetsz, ki ^elviszed az ^üzene^tét.
  \endverse

  \beginverse*
    Refrén \rep{2}
  \endverse

\endsong

\beginsong{A vak ember}[]

  \versesep=12pt plus 3pt minus 8pt

  \ifchorded
    \beginverse*
      {\nolyrics Előjáték: /: \[Am] \[G] \[F] \[E] :/ \rep{2}}
    \endverse
  \fi

  \beginverse\memorize
    Az \[Am]úton a \[G]vak ember \[F]ül és ki\[E]ált, \rep{3}
  \endverse

  \beginchorus
    Oh, oh, \[E7]oh!
    Hol van az {\[Am G F]{\_{ú}}}t? \[E]Hol van a f{\[Am G F]{\_{é}}}ny, \[E]az igazs{\[Am G F]{\_{á}}}g,
    Ami haza\[E]visz? Oh-\[E7]oh-oh, \[E]oh-oh, \[E7]oh-oh!
  \endchorus

  \beginverse
    A ^béna is ^az úton ^ül és ki^ált, \rep{3}
  \endverse

  \beginverse
    A ^néma is ^az úton ^ül és (tátog), \rep{2}
    A \[Am]néma is \[G]az úton \[F]ül és \[E]kiált:
  \endverse

  \beginverse
    S mi \[Am]mind csak ü\[G]lünk, és a \[F]szívünk ki\[E]ált, \rep{3}
  \endverse

  \beginverse
    De \[Am]Jézus is \[G]ott ül ve\[F]lünk és ki\[E]ált, \rep{3}
    Oh, oh, \[E7]oh!
    Én vagyok az {\[Am G F]{\_{ú}}}t, \[E]az igazs{\[Am G F]{\_{á}}}g, és az él{\[Am G F]{\_{e}}}t,
    Ami haza\[E]visz! Oh-\[E7]oh-oh, \[E]oh-oh, \[E7]oh-oh!
  \endverse

  \beginverse*
    \[Am]Jézus \echo{taps} \rep{2}
  \endverse

  \beginverse*
    Oh, oh, \[E7]oh!
    Én vagyok az {\[Am G F]{\_{ú}}}t, \[E]az igazs{\[Am G F]{\_{á}}}g, és az él{\[Am G F]{\_{e}}}t,
    Ami haza\[E]visz! Oh-\[E7]oh-oh, \[E]oh-oh, \[E7]oh-oh!
  \endverse

  \beginverse*
    \[Am]Jézus \echo{taps}
  \endverse

\endsong

\beginsong{Várj és ne félj}
[
  by={Taizé / Jacques Berthier}
]

  \beginverse*
    \[Em]Várj és ne \[C]félj, az \[Am6]Úr jön \[H7]már!
    \[Em]Várj \[D]és ne \[G]félj, hű \[Am7]szív\[H]vel \[Em]várj!
  \endverse

\endsong

\beginsong{A világnak Krisztus kell}
[
  by={Sillye Jenő}
]
  \ifchorded
    \beginverse*
      {\nolyrics Előjáték: \[E] \[H7] \[E] \[A] \[E]}
    \endverse
  \fi

  \beginchorus
    /: A \[E]világnak \[H7] Krisztus \[E]kell!
    A \[A]világnak \[E] Krisztus \[H7]kell!
    A \[E]világnak \[G#]kellesz \[C#m]te is,
    \[A]mivel te \[E]Krisztushoz \[H7]tartoz\[E]ol! :/ \rep{2}
  \endchorus

  \beginverse*
    \[A]Emlé\[E]kezz! \[A]Emlé\[E]kezz, te \[E]Krisztushoz \[H7]tartoz\[E]ol.
  \endverse

\endsong

\input{songs/virrassz-meg.tex}
\beginsong{Vizek felett \\ Oceans}
[
  by={Hillsong United}
]

  \ifchorded
    \beginverse*
      {\nolyrics Előjáték: \[Hm] \[D] \[A] \[G]}
    \endverse
  \fi

  \beginverse\memorize
    \[Hm] Meghívtál, hogy vízre \[D]lépjek, hol nélkü\[A]led elsüllye\[G]dek.
    \[Hm] Ebben megtalállak \[D]Téged, a mélység\[A]ben megt\_{a}rt hi\[G]tem.
  \endverse

  \beginchorus
    \[G] Nagy ne\[D]ved h\_{í}vom \[A]én \[A4] \[G] és felné\[D]zek a vizek fö\[A]lé.
    Ott l\_{á}tlak \[G]én. A lelkem \[D]benned m\_{e}gpi\[A]hen. Enyém v\[G]agy és \[A]én Ti\[Hm]ed.
  \endchorus

  \beginverse
    ^ A mélységnél nagyobb ke^gyelmed, mely elve^zet és tart en^gem.
    ^ Ha elbuknék és nagyon ^félnék, Te nem hagysz ^el és nem inogsz ^meg.
  \endverse

  \beginverse*
    Refrén
  \endverse

  \ifchorded
    \beginverse*
      {\nolyrics Közjáték: /: \[Hm] \[G] \[D] \[A] :/ \rep{2}}
    \endverse
  \fi

  \beginverse*
    /: \[Hm] Lélek add, hogy Benned \[G]teljesen megbízzak,
    A vízen \[D]bátran Veled járjak, és \[A]bárhová hívsz, menjek!
    \[Hm] Vigyél tovább, mint a \[G]lábam tudna menni,
    Taníts \[D]teljes hittel járni, jelen\[A]létedben élni! :/ \rep{3}
  \endverse

  \beginverse*
    Refrén \rep{2} \ifchorded{\nolyrics (Az elsőnek a végén \[D] -vel.)}\fi
  \endverse

  \beginverse*
    Enyém \[G]vagy és \[A]én Ti\[Hm]ed. Enyém \[G]vagy és \[A]én Ti\[D]ed.
  \endverse

\endsong


  \end{songs}

  \songsection{Köszöntések}
  \renewcommand{\thesongnum}{K\arabic{songnum}}
  \begin{songs}{}
    \input{songs-extra.tex}
  \end{songs}

  \songsection{Liturgikus énekek}
  \renewcommand{\thesongnum}{M\arabic{songnum}}
  \setlength{\songnumwidth}{1.25cm}
  \begin{songs}{}
    \input{songs-mass.tex}
  \end{songs}

  \newpage
  \thispagestyle{empty}

  % To be continued...
  \vspace*{\fill}
  % \vfill
  Készítette: Balics Eszter és Donkó István
  % ?Lektorálta: Nobilis Márió?

\end{document}
